\section{Testing}
Il testing è stato realizzato sulla base dell'approccio White Box. Sono stati sviluppati test relativi sia alle principali classi della business logic e della persistenza dei dati, sia ai controller JavaFX, con l'obiettivo di verificare il corretto comportamento delle varie funzionalità.

\subsection{Unit Test}

Sono stati realizzati test unitari per verificare il comportamento delle principali parti del sistema.\\

Per la classe \texttt{Session} è stato verificato il corretto funzionamento del pattern Singleton, controllando che chiamate multiple a getInstance() restituiscano sempre la stessa istanza. È inoltre stata verificata la corretta gestione dell'utente corrente tramite setCurrentUser(), getCurrentUser() e logout().\\
Per la classe \texttt{Database} sono stati verificati l'inserimento degli utenti senza la creazione di duplicati, la corretta gestione delle associazioni tra pazienti e diabetologi, il recupero delle informazioni relative a un paziente a cui sono stati applicati dei filtri nella ricerca, la gestione dei messaggi e le operazioni di inserimento, modifica ed eliminazione di rilevazioni, assunzioni, segnalazioni e terapie. Riguardo alla gestione delle terapie sono inoltre stati verificati casi quali l'assegnazione a un paziente, il comportamento in caso di assegnazione ripetuta e la modifica di una terapia già assegnata; in particolare, che la modifica di una terapia assegnata a più pazienti comporti la modifica per un paziente soltanto.\\
Per \texttt{GestoreAlert} sono stati verificati casi relativi alla mancata assunzione di farmaci, all'aderenza della terapia per più giorni consecutivi e alla regolarità dei valori glicemici prima e dopo i pasti, considerando soglie che attivino l'invio di alert con diversi livelli di urgenza.\\

Sono inoltre stati realizzati test sui controller dell'applicazione: \\
Sono stati inoltre realizzati test sui controller dell'applicazione. Per \texttt{DiabetologoController} sono state verificate l'inizializzazione della schermata, la visualizzazione dei pazienti associati al medico, la ricerca per nome, cognome e codice fiscale, la gestione della ricerca senza risultati e la corretta apertura delle funzionalità relative ad andamento, terapia e informazioni del paziente.

Per la gestione delle terapie sono stati verificati la visualizzazione e la ricerca delle terapie, l'apertura delle finestre per la creazione e la selezione di terapie esistenti, l'assegnazione e la rimozione delle terapie e il comportamento dei relativi pulsanti.

Per \texttt{MessaggiController} sono stati testati la visualizzazione dei messaggi, la ricerca per mittente, la gestione dei messaggi letti e non letti, la visualizzazione del livello di urgenza e l'aggiornamento delle notifiche.

Per \texttt{PazienteController} sono state verificate le operazioni di inserimento e salvataggio di rilevazioni, assunzioni di farmaci, segnalazioni ed email, compresi i casi di compilazione incompleta dei dati.

Per \texttt{StoricoController} sono state verificate la visualizzazione e la ricerca degli elementi dello storico, le operazioni di modifica con e senza salvataggio e le operazioni di eliminazione.

Infine, per \texttt{ResponsabileController}, \texttt{ModificaCredenzialiController} e \texttt{AggiungiPersonaController} sono state verificate la gestione delle liste di utenti, la ricerca, la modifica dei dati, l'associazione tra pazienti e diabetologi e la creazione di nuovi utenti, compresi i controlli sui dati obbligatori e sugli username già presenti.

\subsection{System Test}
 \textcolor{red}{TODO: Tutto}\\
(Test sull'applicazione completa)\\
Probabili casi da verificare:\\
-Registrazione di una glicemia fuori soglia -$>$ alert $>$ nella dashboard del diabetologo di riferimento\\
-Mancata registrazione di assunzioni farmaco per più di 3 giorni consecutivi -$>$ alert generato al medico\\
-Modifica di una terapia da parte del medico -$>$ le assunzioni successive vengono validate sulla nuova terapia\\
-Login con credenziali errate -$>$ accesso rifiutato\\
-Impossibilità di accedere a dati altrui per chi non ne ha le facoltà\\


